\clearpage\SourceLeaf{186}%
\pdfbookmark[1]{Neue Aussicht zur Erweiterung des Gebiets der Himmelskunde}{aussicht1813}
\fancyhead[L]{\small NEUE AUSSICHT ZUR ERWEITERUNG DER HIMMELSKUNDE}
\begin{center}
{\small NEUE AUSSICHT\par
ZUR ERWEITERUNG DES GEBIETS DER HIMMELSKUNDE.\par}
\medskip\hrule\medskip
{\small Neues Göttingisches Taschenbuch zum Nutzen und Vergnügen für das Jahr 1813.\par
Göttingen, bei Heinrich Dieterich.\par}
\medskip\hrule
\end{center}\medskip
So weit die Geschichte reicht, war die Erde stets der unruhige Tummelplatz feindseliger Leidenschaften. Eine Völkerschaft verdrängt die andere, um früh oder spät wieder einer andern Platz zu machen; ein Reich erhebt sich auf den Trümmern des andern: Vergänglichkeit ist das Loos aller menschlichen Schöpfungen, deren Segen oder Fluch die Generationen, welche eben der Strom der Zeit an ihnen vorüber führt, theilen. Demüthigend und niederschlagend ist dieser ewige Wechsel aller Dinge auf der Erde, aber erhebend und erheiternd ist ein Blick auf das Schauspiel des steten Fortschreitens des menschlichen Geistes im Gebiete der Wahrheit, und in keinem Theile desselben mehr, als in der erhabensten der Naturwissenschaften, der Astronomie.

Seit zwei Jahrtausenden sehen wir die schöne Frucht des Beobachtens und des tiefen Nachdenkens, mitten in den Stürmen der Zeiten, von liebenden Händen der Geweihten gepflegt, sich entwickeln und wachsen; ein Jahrhundert übergibt dem andern das kostbare Vermächtniss, es weiter zu mehren. Immer mehr entschleiert die Himmlische von ihren wundervollen Geheimnissen. Gezählt sind die Sterne; gemessen ist die unermesslich scheinende Weite der Sonne und der sie umkreisenden Weltkörper; genau bestimmt sind ihre Bahnen: vollständig erklärt, aus Einer erhaben einfachen Grundkraft, alle ihre verwickelten Bewegungen und alle mannigfaltigen daraus hervorgehenden Erscheinungen. Gemessen,
\expandafter\def\csname HistoricalNoteText13\endcsname{\textsc{Piazzi}'s, \textsc{Calandrelli}'s und \textsc{Brinkley}'s Beobachtungen geben zwar eine kleine jährliche Parallaxe bei dem hellen Sterne in der Leier, allein die Beobachtungen sind noch Zweifeln unterworfen, und die Resultate stimmen nicht überein; auch folgt bei demselben Sterne aus \textsc{Bradley}'s Beobachtungen gar keine merkliche Parallaxe.}%
\expandafter\def\csname HistoricalNoteText14\endcsname{Vermittelst der Durchgänge von 1761 und 1769.}%
\expandafter\def\csname HistoricalNoteText15\endcsname{Mit Hülfe der einen von der Sonnenparallaxe abhangenden kleinen Ungleichheit in der Bewegung des Mondes.}%
\SourceLeaf{187}%
ja gewogen, sind Erde, Sonne, Planeten und Monden; gemessen ist die Schnelligkeit des Lichtstrahls. Verschwunden sind die Wahnbegriffe, die ehedem die Völker ängstigten. Auch dem Irdischen zu dienen, verschmäht die Himmlische nicht. Sie theilt und zählt dem Menschen seine Zeiträume, misst und zeichnet ihm Länder und Welttheile, und leitet seine Schiffe sicher über den Ocean.

Unser Sonnensystem ist der eigentliche Schauplatz unsrer bisherigen Astronomie: noch viele Jahrhunderte hindurch wird es unerschöpften Stoff zu neuen Forschungen geben. Aber noch viel weniger, als die Erde in unserm Sonnensystem, ist dieses im Weltall: ein Tropfen im Ocean. Das Heer von Fixsternen ist für uns bis jetzt nicht viel mehr, als eine Zahl fester Punkte, an die wir unsre Beobachtungen anknüpfen. Nur eine Ahnung haben wir von ihren unermesslichen Fernen: fast ganz verschwindet gegen diese die grösste Grundlinie, mit welcher wir sie zu messen unternehmen, der Durchmesser der jährlichen Erdbahn\HistoricalNote{13}. Geben wir indess darum die Hoffnung, von diesen Weiten eine Kenntniss zu erhalten, nicht auf. Auch die Weite der Sonne hat man mit dem kleinen Erddurchmesser nicht unmittelbar gemessen, sondern statt jener unter günstigen Umständen, die nähere Venus gewählt{\renewcommand{\thefootnote}{**)}\HistoricalNote{14}}, und später hat der Scharfsinn die schon länger bekannte Entfernung des Mondes zu einem Maassstabe für die Entfernung der Sonne zu gebrauchen gewusst{\renewcommand{\thefootnote}{***)}\HistoricalNote{15}}. Vielleicht sind gerade die hellsten Fixsterne uns nicht die allernächsten: was bei jenen bisher noch nicht ganz hat gelingen wollen, glückt vielleicht besser bei einem unscheinbaren Stern von einer viel niedrigern Ordnung.

In der That gibt es für diese Vermuthung wichtige Gründe. Die sogenannten Fixsterne, zu welchen auch unsere Sonne gehört, sind im strengen Sinne des Worts nicht fest; wären sie auch einmal alle zugleich in Ruhe gewesen, so hätten sie es doch wegen ihrer gegenseitigen Einwirkungen nicht bleiben können. Auch in dem zahllosen Fixsternheere ist überall Bewegung, aber wie hier die Räume mit grossen Maassstäben gemessen werden müssen, so auch die Zeiten.
\SourceLeaf{188}%
Ein Jahrhundert bringt in ihre gegenseitigen Stellungen keine so grossen Änderungen, als ein Tag in die gegenseitigen Stellungen der Planeten. Aber merklich sind unsern heutigen feinen Beobachtungen diese veränderten Stellungen doch schon seit funfzig oder hundert Jahren bei einer beträchtlichen Anzahl von Sternen. Unter gleichen Umständen werden diese Änderungen bei den \emph{nächsten} Sternen am meisten in die Augen fallen, und sonderbar genug, sind sie bei einigen kleinen Sternen beträchtlicher, als bei den Sternen erster Grösse. \emph{Solche} Sterne, von denen man vielleicht noch mehrere entdecken wird, sollte man vorzüglich in Beziehung auf ihre jährliche Parallaxe, d.\ i.\ ihre verschiedene Stellung nach dem verschiedenen Stande der Erde in ihrer Bahn, recht sorgfältig beobachten.

Besonders ist ein Stern in dieser Rücksicht der höchsten Aufmerksamkeit würdig. Es ist derjenige im Sternbilde des Schwans, von der sechsten Grösse, welchen \textsc{Flamstead} mit der Zahl 61 bezeichnet hat. Unter allen Sternen, an welchen man bisher eine Verrückung wahrgenommen hat, zeigt dieser die stärkste, nemlich 8 Min.\ 52 Sec.\ in hundert Jahren. Eine so beträchtliche Veränderung der Stellung in Beziehung auf uns, setzt entweder ein sehr schnelles wirkliches Fortrücken im Weltraume voraus, oder eine verhältnissmässig geringere Entfernung. Wir verdanken jene wichtige Bemerkung Herrn \textsc{Bessel} in Königsberg, welchen die Vergleichung von \textsc{Bradley}'s Beobachtungen dieses Sterns mit Beobachtungen aus den letzten Jahren des vorigen Jahrhunderts darauf geführt hat.

Allein noch weit interessanter wird die beobachtete schnelle Bewegung dieses Sterns durch einen andern Umstand. Der Stern ist eigentlich ein sogenannter \emph{Doppelstern}, das heisst, durch gute Fernröhre sieht man, dass er aus zweien sehr nahe bei einander erscheinenden Sternen besteht (ihre Entfernung von einander beträgt 20 Secunden, und beide Sterne sind an Helligkeit nicht sehr verschieden). Dergleichen Doppelsterne gibt es am Himmel in grosser Menge, und die Astronomen waren bisher immer der Meinung, dass zwei Sterne uns nur zufällig als Doppelsterne erscheinen, indem sie mit uns beinahe in Einer geraden Linie liegen, wobei also ihre Entfernungen sehr ungleich sein können. Vor wenigen Jahren hat zuerst Hr.\ \textsc{Herschel} diese Meinung bestritten, und behauptet, die Doppelsterne stehen wirklich im Raume nahe bei einander und machen ein besonderes System aus. \textsc{Herschels} Gründe für diese Meinung sind freilich nicht entscheidend, aber an sich schon ist es in einem äusserst hohen Grade unwahr%
\SourceLeaf{189}%
scheinlich, dass der Zufall so viele Doppelsterne mit so sehr kleinem Zwischenraume gebildet haben sollte; es lässt sich viel eher für eine wirkliche Beziehung zwischen denselben wetten. Unser Stern im Schwan gibt nunmehr einen grossen Ausschlag für \textsc{Herschels} Meinung. Stehen beide Sterne in keiner weitern Verbindung mit einander, und verändern sie ihren Platz im Weltraume oder wir den unsrigen, so lässt sich mit höchster Wahrscheinlichkeit erwarten, dass sie bald aufhören werden, uns als ein Doppelstern zu erscheinen. So zeigen sich öfters am Himmel nahe beim Jupiter kleine Sterne; aber wie der Planet fortrückt, scheint er die Sterne zu verlassen, nur die ihm wirklich nahen Trabanten begleiten ihn beständig treu. Unser Doppelstern ist aber in 50 Jahren dreizehnmal so weit am Firmament fortgerückt, als der Zwischenraum zwischen beiden Sternen beträgt, und der Zwischenraum selbst ist fast unverändert geblieben, nur weniger schief stehen die Sterne jetzt gegen einander (im Jahre 1755 war der eine nordöstlich vom andern, jetzt fast ganz östlich). Es bleibt also kaum ein Zweifel, dass wir in ihnen eine neue Merkwürdigkeit des Weltbaus kennen gelernt haben, ein Sonnensystem, worin nicht Eine, sondern zwei Sonnen herrschen.

Ein blosses reines Herrschen findet in der Körperwelt nicht Statt, das Herrschende empfindet jedesmal mehr oder weniger die Rückwirkung des Beherrschten. Jene beiden Sonnen werden eine die andere wechselseitig beherrschen, die zweite wird nicht um die erste, die erste nicht um die zweite laufen, sondern beide um ihren gemeinschaftlichen Schwerpunkt. Ein solcher Wechselumlauf wird desto schneller vollendet werden, je näher beide einander sind, und je kräftiger sie auf einander einwirken. Wie weit sie von einander sind, werden wir beurtheilen können, wenn die Hoffnung gegründet ist, ihre Entfernung von uns noch messbar zu finden. Und die Dauer jenes Umlaufs, der vielleicht nur wenige Jahrhunderte beträgt, wird die Zeit schon lehren. Dann wird man einen grossen Schritt in der Kenntniss des Sternenhimmels weiter gekommen sein, man wird dann im Stande sein, jene Sterne \emph{abzuwägen}, wenn auch nicht sie einzeln, doch ihre Gesammtmasse. Der erste Schritt wird nicht der einzige bleiben. Aber wer zählt die Jahrtausende, die erst verfliessen müssen, bis unsere Kenntnisse vom Sternhimmel nur ein Schatten von denen werden können, die wir längst von \emph{unserm} so grossen und so kleinen Sonnensystem besitzen, und wie viel bleibt nicht noch in diesem immer Neues zu lernen übrig.
